\section{Kiến trúc hệ thống}

\subsection{Sơ đồ kiến trúc}

\begin{figure}[H]
    \centering
    \fbox{\parbox{0.9\textwidth}{\centering TODO: Chèn sơ đồ kiến trúc hệ thống tại đây. Bản ASCII tham khảo nằm trong file \texttt{ARCHITECTURE\_ASCII.txt}.}}
    \caption{Sơ đồ kiến trúc serverless của ứng dụng Task Manager}
\end{figure}

\subsection{Các thành phần chính}

\begin{table}[H]
\centering
\begin{tabular}{|p{4cm}|p{9cm}|}
\hline
\textbf{Thành phần} & \textbf{Vai trò} \\
\hline
Amazon S3 & Lưu trữ các file frontend. Bucket được đặt private, bật toàn bộ Block Public Access và không bật Static Website Hosting. \\
\hline
Amazon CloudFront & CDN duy nhất phục vụ frontend cho người dùng. CloudFront lấy file từ S3 thông qua OAC. \\
\hline
Origin Access Control & Cho phép CloudFront truy cập S3 private bucket một cách an toàn, không cần public bucket. \\
\hline
Amazon Cognito & Quản lý đăng ký, đăng nhập và phát hành JWT token cho frontend. \\
\hline
API Gateway REST API & Cung cấp HTTPS endpoint, định tuyến request đến Lambda, kiểm tra JWT bằng Cognito Authorizer và áp dụng throttling. \\
\hline
AWS Lambda & Chạy backend serverless bằng Node.js 20.x. Bốn Lambda function tách riêng tương ứng với bốn thao tác CRUD. \\
\hline
Amazon VPC & Cô lập tầng compute trong custom VPC \texttt{TaskManagerVPC} với hai private subnet. \\
\hline
DynamoDB Gateway Endpoint & Định tuyến traffic từ Lambda đến DynamoDB qua AWS private backbone, không cần NAT Gateway. \\
\hline
Amazon DynamoDB & Lưu trữ dữ liệu công việc trong bảng \texttt{TasksTable}, có GSI \texttt{userId-index}. \\
\hline
IAM & Cấp quyền least privilege bằng các role riêng cho từng Lambda function. \\
\hline
CloudWatch, SNS, AWS Budget & Giám sát metric, log, cảnh báo lỗi qua email và kiểm soát chi phí hàng tháng. \\
\hline
\end{tabular}
\caption{Vai trò của các dịch vụ AWS trong hệ thống}
\end{table}

\subsection{Resource chính}

\begin{table}[H]
\centering
\begin{tabular}{|p{5cm}|p{8cm}|}
\hline
\textbf{Loại resource} & \textbf{Tên đã sử dụng} \\
\hline
VPC & \texttt{TaskManagerVPC} \\
\hline
Private subnet AZ-1 & \texttt{PrivateSubnet-AZ1}, CIDR \texttt{10.0.1.0/24}, AZ \texttt{ap-southeast-1a} \\
\hline
Private subnet AZ-2 & \texttt{PrivateSubnet-AZ2}, CIDR \texttt{10.0.2.0/24}, AZ \texttt{ap-southeast-1b} \\
\hline
Security group & \texttt{LambdaSecurityGroup} \\
\hline
DynamoDB table & \texttt{TasksTable} \\
\hline
DynamoDB GSI & \texttt{userId-index} \\
\hline
Cognito User Pool & \texttt{TaskManagerUserPool} \\
\hline
Cognito App Client & \texttt{WebAppClient} \\
\hline
API Gateway REST API & \texttt{TaskManagerAPI} \\
\hline
CloudWatch Dashboard & \texttt{TaskManager-Dashboard} \\
\hline
AWS Budget & \texttt{TaskManager-Budget-0.01} \\
\hline
\end{tabular}
\caption{Các resource chính của hệ thống}
\end{table}

\subsection{Luồng request}

Khi người dùng mở ứng dụng, trình duyệt truy cập CloudFront domain. Nếu các file tĩnh như \texttt{index.html}, CSS và JavaScript đã có trong cache, CloudFront trả về trực tiếp từ edge location. Nếu cache miss, CloudFront lấy object từ S3 private bucket thông qua OAC. S3 bucket không public nên truy cập trực tiếp đến URL S3 sẽ bị từ chối.

Luồng tạo công việc sau khi người dùng đăng nhập:

\begin{enumerate}
    \item Người dùng đăng nhập bằng Cognito và nhận JWT token.
    \item Frontend gửi request HTTPS đến API Gateway endpoint \texttt{/tasks}, kèm header \texttt{Authorization}.
    \item API Gateway dùng Cognito Authorizer để xác thực JWT token.
    \item Nếu token hợp lệ, API Gateway gọi Lambda function tương ứng.
    \item Lambda function chạy trong private subnet của \texttt{TaskManagerVPC}.
    \item Lambda gửi request đến DynamoDB qua route table có destination là DynamoDB prefix list và target là VPC endpoint.
    \item DynamoDB ghi hoặc đọc dữ liệu trong \texttt{TasksTable}.
    \item Lambda trả JSON response về API Gateway, sau đó API Gateway trả kết quả về trình duyệt.
\end{enumerate}

\subsection{Tự mở rộng serverless}

Hệ thống không cần Application Load Balancer vì API Gateway đã đóng vai trò entry point HTTPS và định tuyến request đến Lambda. Khi có nhiều request đồng thời, Lambda tự động tạo thêm execution environment để xử lý song song. Khi không có request, Lambda có thể giảm về không instance đang chạy, giúp giảm chi phí vận hành.

Để tránh runaway cost và bảo vệ giới hạn Free Tier, reserved concurrency được cấu hình như sau:

\begin{table}[H]
\centering
\begin{tabular}{|p{5cm}|p{3cm}|p{5cm}|}
\hline
\textbf{Lambda function} & \textbf{Reserved concurrency} & \textbf{Lý do} \\
\hline
\texttt{GetTasksFunction} & 5 & Cho phép nhiều request đọc danh sách song song trong demo. \\
\hline
\texttt{CreateTaskFunction} & 2 & Giới hạn thao tác ghi mới để kiểm soát chi phí và lỗi lặp request. \\
\hline
\texttt{UpdateTaskFunction} & 2 & Đủ cho demo cập nhật đồng thời, vẫn giới hạn mức ghi. \\
\hline
\texttt{DeleteTaskFunction} & 1 & Xóa task là thao tác ít dùng hơn và nên được giới hạn chặt. \\
\hline
\end{tabular}
\caption{Reserved concurrency của các Lambda function}
\end{table}

Nếu reserved concurrency đặt quá thấp, người dùng có thể gặp throttling dù hệ thống vẫn còn khả năng xử lý. Nếu đặt quá cao, lỗi vòng lặp hoặc lượng request bất thường có thể tạo nhiều invocation và làm tăng chi phí.

\subsection{CDN và Origin Access Control}

CloudFront giúp giảm độ trễ bằng cách cache static assets tại edge location gần người dùng. Khi cache hit, CloudFront trả nội dung mà không cần truy cập lại S3. Khi cache miss, CloudFront lấy nội dung từ origin là S3 bucket.

S3 bucket bắt buộc private vì frontend không được truy cập qua URL S3 trực tiếp. Origin Access Control là cơ chế cho phép CloudFront ký request đến S3 và bucket policy chỉ cho phép CloudFront service principal đọc object từ distribution hợp lệ. Nhờ đó, CloudFront là public entry point duy nhất của frontend.

\subsection{VPC Endpoint}

Lambda được gắn VpcConfig vào hai private subnet để backend chạy trong mạng riêng. Vì private subnet không có NAT Gateway, Lambda không thể gọi DynamoDB qua internet công cộng. DynamoDB Gateway Endpoint giải quyết vấn đề này bằng cách thêm route đến DynamoDB prefix list trong route table của private subnet.

Traffic từ Lambda đến DynamoDB đi qua AWS private backbone theo đường: Lambda trong private subnet $\rightarrow$ route table $\rightarrow$ DynamoDB Gateway Endpoint $\rightarrow$ DynamoDB. Cách này an toàn hơn và không phát sinh chi phí cố định như NAT Gateway.
